\section*{Chapitre \ref{ch:g12:seq} --- Suites}

\begin{solution}{exo:g12:seq:1}
Soit $P(n)$ l'énoncé $\sum_{k=1}^{n} k^2 = \frac{n(n+1)(2n+1)}{6}$.
\emph{Initialisation~:} pour $n = 0$ les deux membres valent $0$ (somme
vide).
\emph{Hérédité~:} supposons $P(n)$. Alors
\[
\sum_{k=1}^{n+1} k^2 = \frac{n(n+1)(2n+1)}{6} + (n+1)^2
= \frac{(n+1)\bigl(n(2n+1) + 6(n+1)\bigr)}{6}
= \frac{(n+1)(2n^2 + 7n + 6)}{6}.
\]
Comme $2n^2 + 7n + 6 = (n+2)(2n+3)$, c'est
$\frac{(n+1)(n+2)(2(n+1)+1)}{6}$, qui est $P(n+1)$. Par récurrence,
$P(n)$ vaut pour tout $n$.
\end{solution}

\begin{solution}{exo:g12:seq:2}
$a_{n+1} - a_n = \frac{n+2}{n+1} - \frac{n+1}{n}
= \frac{n(n+2) - (n+1)^2}{n(n+1)} = \frac{-1}{n(n+1)} < 0$~: $(a_n)$ est
strictement \omterm{def:g10:functions:variations}{décroissante}.

$(b_n)$ a des termes positifs et
$\frac{b_{n+1}}{b_n} = \frac{2^{n+1}}{n+1}\cdot\frac{n}{2^n} =
\frac{2n}{n+1} \geq 1 \iff 2n \geq n+1 \iff n \geq 1$~: $(b_n)$ est
\omterm{def:g12:seq:monotonic}{croissante} (strictement pour $n \geq 2$).

$c_{n+1} - c_n = (n+1)^2 - 10(n+1) - n^2 + 10n = 2n - 9$, qui est négatif
pour $n \leq 4$ et positif pour $n \geq 5$~: $(c_n)$ décroît jusqu'à
$c_5 = -25$, son \omterm{def:g10:functions:extrema}{minimum}, puis croît. Elle n'est pas \omterm{def:g11:func:monotone}{monotone}.
\end{solution}

\begin{solution}{exo:g12:seq:3}
Factoriser les termes dominants~:
\[
u_n = \frac{n^2(3 - 1/n + 1/n^2)}{n^2(2 + 5/n^2)} \xrightarrow[n\to+\infty]{} \frac{3}{2}.
\]
Multiplier par le conjugué~:
\[
v_n = \frac{(n+1) - n}{\sqrt{n+1} + \sqrt{n}} = \frac{1}{\sqrt{n+1}+\sqrt{n}}
\xrightarrow[n\to+\infty]{} 0.
\]
Diviser numérateur et dénominateur par $3^n$~:
\[
w_n = \frac{(2/3)^n - 1}{1 + (1/3)^n} \xrightarrow[n\to+\infty]{} \frac{0-1}{1+0} = -1,
\]
en utilisant $\lim q^n = 0$ pour $\abs{q} < 1$.
\end{solution}

\begin{solution}{exo:g12:seq:4}
$u_n = 5 + 3n \to +\infty$ et $v_n = 8 \cdot (1/2)^n = 2^{3-n} \to 0$.
Les sommes sont
\[
\sum_{k=0}^{n} u_k = (n+1)\,\frac{5 + (5+3n)}{2} = \frac{(n+1)(10+3n)}{2}
\xrightarrow[n\to+\infty]{} +\infty,
\]
\[
\sum_{k=0}^{n} v_k = 8\,\frac{1 - (1/2)^{n+1}}{1 - 1/2}
= 16\left(1 - \left(\tfrac{1}{2}\right)^{n+1}\right)
\xrightarrow[n\to+\infty]{} 16 .
\]
\end{solution}

\begin{solution}{exo:g12:seq:5}
Comme $-1 \leq \cos n \leq 1$,
\[
\frac{n-1}{n+1} \leq \frac{n + \cos n}{n+1} \leq 1,
\]
et $\frac{n-1}{n+1} \to 1$, donc la limite est $1$ par le théorème des
gendarmes.

Pour la seconde limite, écrire
\[
0 \leq \frac{n!}{n^n}
= \frac{1}{n}\cdot\frac{2}{n}\cdots\frac{n}{n}
\leq \frac{1}{n},
\]
car chaque facteur $\frac{k}{n}$ avec $2 \leq k \leq n$ est au plus $1$.
Comme $\frac1n \to 0$, le théorème des gendarmes donne
$\frac{n!}{n^n} \to 0$. (La borne \omterm{def:g12:seq:arith-geom}{géométrique} suggérée fonctionne aussi~:
chaque facteur avec $k \leq n/2$ est au plus $\frac12$, donnant la borne
plus forte $(1/2)^{\floor{n/2}}$.)
\end{solution}

\begin{solution}{exo:g12:seq:6}
\emph{1.} $u_0 = 0 \in \intcc{0}{2}$. Si $0 \leq u_n \leq 2$, alors
$2 \leq u_n + 2 \leq 4$, donc $\sqrt{2} \leq u_{n+1} \leq 2$~; en
particulier $0 \leq u_{n+1} \leq 2$. Par récurrence la propriété vaut pour
tout $n$.

\emph{2.} $u_{n+1} - u_n = \sqrt{u_n + 2} - u_n$. Pour
$x \in \intcc{0}{2}$, $\sqrt{x+2} \geq x \iff x + 2 \geq x^2 \iff
(2-x)(x+1) \geq 0$, ce qui est vrai. D'où $(u_n)$ est \omterm{def:g12:seq:monotonic}{croissante}.

\emph{3.} \omterm{def:g12:seq:monotonic}{Croissante} et \omterm{def:g12:seq:bounded}{majorée} par $2$, $(u_n)$ \omterm{def:g12:seq:limit}{converge} vers un certain
$\ell \in \intcc{0}{2}$. En passant à la limite dans
$u_{n+1} = \sqrt{u_n + 2}$ (l'application $x \mapsto \sqrt{x+2}$ est
continue) on obtient $\ell = \sqrt{\ell + 2}$, donc
$\ell^2 - \ell - 2 = 0$, c'est-à-dire $\ell \in \{-1, 2\}$. Comme
$\ell \geq 0$, $\lim u_n = 2$.
\end{solution}

\begin{solution}{exo:g12:seq:7}
\emph{1.} Entre deux doses, $40\,\%$ du médicament est éliminé, donc la
quantité $u_n$ devient $0.6\,u_n$~; la dose suivante ajoute $1$ unité~:
$u_{n+1} = 0.6\,u_n + 1$.

\emph{2.} $v_{n+1} = u_{n+1} - 2.5 = 0.6\,u_n + 1 - 2.5 = 0.6(u_n - 2.5)
= 0.6\,v_n$~: $(v_n)$ est \omterm{def:g12:seq:arith-geom}{géométrique} de raison $0.6$ et de premier terme
$v_0 = 1 - 2.5 = -1.5$. D'où $v_n = -1.5 \times 0.6^n$ et
\[
u_n = 2.5 - 1.5 \times 0.6^n .
\]

\emph{3.} Comme $0.6^n \to 0$, $u_n \to 2.5$~: la quantité de médicament
se stabilise à $2.5$ unités.
\end{solution}

\begin{solution}{exo:g12:seq:8}
\emph{1.} $u_0 = 3 > 1$. Si $u_n > 1$, alors $u_n + 2 > 0$ et
\[
u_{n+1} - 1 = \frac{4u_n - 1 - u_n - 2}{u_n + 2} = \frac{3(u_n - 1)}{u_n + 2} > 0 .
\]
Par récurrence, $u_n > 1$ pour tout $n$ (et en particulier
$u_n + 2 \neq 0$, donc la \omterm{def:g12:seq:sequence}{suite} est bien définie).

\emph{2.} En utilisant l'identité ci-dessus,
\[
v_{n+1} = \frac{1}{u_{n+1} - 1} = \frac{u_n + 2}{3(u_n - 1)}
= \frac{(u_n - 1) + 3}{3(u_n - 1)} = \frac{1}{3} + v_n .
\]
Donc $(v_n)$ est \omterm{def:g12:seq:arith-geom}{arithmétique} de raison $\frac13$ et
$v_0 = \frac{1}{u_0 - 1} = \frac12$.

\emph{3.} $v_n = \frac12 + \frac{n}{3}$, d'où
$u_n = 1 + \frac{1}{v_n} = 1 + \frac{6}{3 + 2n}$. Comme
$v_n \to +\infty$, $u_n \to 1$.
\end{solution}

\begin{solution}{exo:g12:seq:9}
\emph{1.} $H_{2n} - H_n = \sum_{k=n+1}^{2n} \frac{1}{k}$ est une somme de
$n$ termes, chacun au moins $\frac{1}{2n}$~; d'où
$H_{2n} - H_n \geq n \cdot \frac{1}{2n} = \frac12$.

\emph{2.} Par récurrence sur $k$~: $H_{2^0} = H_1 = 1 \geq 1$. Si
$H_{2^k} \geq 1 + \frac{k}{2}$, alors en appliquant le point~1 avec
$n = 2^k$,
\[
H_{2^{k+1}} \geq H_{2^k} + \frac12 \geq 1 + \frac{k+1}{2}.
\]
La \omterm{def:g12:seq:sequence}{suite} $(H_n)$ est \omterm{def:g12:seq:monotonic}{croissante} (chaque pas ajoute $\frac{1}{n+1} > 0$) et
la \omterm{def:g12:seq:sequence}{sous-suite} $H_{2^k}$ est non \omterm{def:g12:seq:bounded}{majorée}, donc $(H_n)$ n'est pas \omterm{def:g12:seq:bounded}{majorée}.
\omterm{def:g12:seq:monotonic}{Croissante} et non \omterm{def:g12:seq:bounded}{majorée}, elle tend vers $+\infty$
(\cref{thm:g12:seq:monotone}).
\end{solution}

\begin{solution}{exo:g12:seq:10}
\emph{1.} La \omterm{def:g12:seq:sequence}{suite} $d_n = b_n - a_n$ satisfait
$d_{n+1} - d_n = (b_{n+1} - b_n) - (a_{n+1} - a_n) \leq 0$, donc $(d_n)$
est \omterm{def:g10:functions:variations}{décroissante}~; comme $d_n \to 0$, on obtient $d_n \geq 0$ pour tout
$n$ (une \omterm{def:g12:seq:sequence}{suite} \omterm{def:g10:functions:variations}{décroissante} avec un terme négatif resterait en dessous
pour toujours, empêchant la limite $0$). D'où $a_n \leq b_n$.

\emph{2.} De $a_n \leq b_n \leq b_0$, la \omterm{def:g12:seq:sequence}{suite} \omterm{def:g12:seq:monotonic}{croissante} $(a_n)$ est
\omterm{def:g12:seq:bounded}{majorée}, donc elle \omterm{def:g12:seq:limit}{converge} vers un certain $\ell$. De même $(b_n)$,
\omterm{def:g10:functions:variations}{décroissante} et minorée par $a_0$, \omterm{def:g12:seq:limit}{converge} vers un certain $\ell'$.
Alors $\ell' - \ell = \lim (b_n - a_n) = 0$, donc $\ell = \ell'$.

\emph{3.} $(a_n)$ est (strictement) \omterm{def:g12:seq:monotonic}{croissante} car
$a_{n+1} - a_n = \frac{1}{(n+1)!} > 0$. Pour $(b_n)$,
\[
b_{n+1} - b_n = \frac{1}{(n+1)!} + \frac{1}{(n+1)(n+1)!} - \frac{1}{n\,n!}
= \frac{n(n+1) + n - (n+1)^2}{n(n+1)(n+1)!}
= \frac{-1}{n(n+1)(n+1)!} < 0 ,
\]
donc $(b_n)$ est \omterm{def:g10:functions:variations}{décroissante}. Enfin $b_n - a_n = \frac{1}{n\,n!} \to 0$.
Les deux \omterm{def:g12:seq:sequence}{suites} sont adjacentes, donc convergent vers une limite
commune.
\end{solution}

\begin{solution}{pb:g12:seq:1}
\textbf{1.} Vrai pour $n = 1$ ($1 = 1^2$). Si
$1 + 3 + \dots + (2n - 1) = n^2$, alors en ajoutant l'impair
suivant~: $n^2 + (2n + 1) = (n + 1)^2$~: hérédité. Par récurrence,
c'est vrai pour tout $n \geq 1$.

\textbf{2.} $2^0 = 1 > 0$. Si $2^n > n$, alors
$2^{n+1} = 2 \cdot 2^n > 2n \geq n + 1$ pour $n \geq 1$ (et le cas
$n = 0$ se vérifie directement)~: hérédité, c'est fait.

\textbf{3.} $n = 0$~: $1 \geq 1$. Si $(1 + x)^n \geq 1 + nx$,
multiplions par $1 + x \geq 1 > 0$~:
$(1 + x)^{n+1} \geq (1 + nx)(1 + x) = 1 + (n + 1)x + nx^2
\geq 1 + (n + 1)x$.

\textbf{4.} C'est le passage de $n = 1$ à $n = 2$~: parmi deux
billes, «~les $n$ premières~» et «~les $n$ dernières~» sont deux
billes uniques \emph{disjointes} --- aucune bille commune ne fait le
pont entre les deux groupes, donc rien n'oblige leurs couleurs à
coïncider. L'argument d'hérédité exige en silence que les deux
groupes se chevauchent, ce qui n'est vrai qu'à partir de
$n \geq 2$~; avec l'initialisation en $n = 1$, la chaîne ne démarre
jamais.

\textbf{5.} $4^0 - 1 = 0 = 3 \times 0$. Si $4^n - 1 = 3k$, alors
$4^{n+1} - 1 = 4(4^n - 1) + 3 = 3(4k + 1)$~: hérédité.

\textbf{6.} $x_1 = \frac32$, $x_2 = \frac{17}{12}$,
$x_3 = \frac{577}{408}$.

\textbf{7.} $x_{n+1}^2 - 2 = \frac{(x_n^2 + 2)^2 - 8x_n^2}
{4x_n^2} = \frac{(x_n^2 - 2)^2}{4 x_n^2}$~: un carré divisé par un
nombre positif, donc $\geq 0$, et même $> 0$ dès que
$x_n^2 \neq 2$. Récurrence~: $x_0 = 2 > 0$ avec $x_0^2 = 4 > 2$~;
si $x_n > 0$ et $x_n^2 > 2$, alors $x_{n+1}$ (\omterm{def:g11:stat:mean}{moyenne} de nombres
positifs) est positif et $x_{n+1}^2 - 2 > 0$.

\textbf{8.} $x_{n+1} - x_n = \frac{2 - x_n^2}{2x_n} < 0$ d'après la
question 7~: la \omterm{def:g12:seq:sequence}{suite} est strictement \omterm{def:g10:functions:variations}{décroissante}.

\textbf{9.} \omterm{def:g10:functions:variations}{Décroissante} et minorée (par $1$, puisque $x_n^2 > 2 >
1$ et $x_n > 0$)~: par le théorème de convergence \omterm{def:g11:func:monotone}{monotone},
$(x_n)$ \omterm{def:g12:seq:limit}{converge} vers une limite $L \geq 1$.

\textbf{10.} Les limites respectent l'algèbre~: de
$x_{n+1} = \frac12\left(x_n + \frac{2}{x_n}\right)$ et
$x_n \to L \geq 1 > 0$ on tire
$L = \frac12\left(L + \frac2L\right)$, donc $L^2 = 2$ et, $L$ étant
positif, $L = \sqrt2$. Verdict~: convergence démontrée, limite
identifiée --- Héron est acquitté avec les honneurs.

\textbf{11.} $x_{n+1} - \sqrt2 = \frac{x_n^2 - 2\sqrt2\,x_n +
2}{2x_n} = \frac{(x_n - \sqrt2)^2}{2x_n}$~: c'est exactement
$e_{n+1} = \frac{e_n^2}{2x_n}$, et $x_n > \sqrt2$ donne
$e_{n+1} \leq \frac{e_n^2}{2\sqrt2}$. Erreur élevée au carré~:
chaque pas double le nombre de décimales exactes, comme on
l'observait déjà dans le volume précédent.

\textbf{12.} $e_0 \approx 0.5858$, $e_1 \approx 0.0858$,
$e_2 \approx 0.00245$, $e_3 \approx 2.1 \times 10^{-6}$.
Rapports~: $\frac{e_1}{e_0^2} \approx 0.25 = \frac{1}{2x_0}$~;
$\frac{e_2}{e_1^2} \approx 0.333 = \frac{1}{2x_1}$~;
$\frac{e_3}{e_2^2} \approx 0.353 \approx \frac{1}{2x_2}$~: le
théorème à l'œuvre.

\textbf{13.} $H_{2^{18}} \geq 1 + 9 = 10$~: il faut environ
$260\,000$ termes ($2^{18} = 262\,144$) rien que pour dépasser $10$
--- une divergence au pas de tortue (et $H_n > 100$ demanderait plus
de termes qu'il n'y a d'atomes dans une bibliothèque).

\textbf{14.} $S_n = 2 - \frac{1}{2^n}$ (somme \omterm{def:g12:seq:arith-geom}{géométrique}), et
$\frac{1}{2^n} \to 0$ (\cref{thm:g12:seq:geometric}), donc
$S_n \to 2$~: la tablette de chocolat indéfiniment croquée tend
vers le tout sans jamais l'atteindre --- dans la langue officielle
des limites, cette fois.

\textbf{15.} Pour $k \geq 2$~:
$\frac{1}{k^2} \leq \frac{1}{k(k-1)} = \frac{1}{k-1} - \frac1k$,
donc $S_n \leq 1 + \left(1 - \frac1n\right) < 2$~: la \omterm{def:g12:seq:sequence}{suite} est
\omterm{def:g12:seq:monotonic}{croissante} et \omterm{def:g12:seq:bounded}{majorée}, donc convergente (convergence \omterm{def:g11:func:monotone}{monotone}).
Euler a plus tard nommé la limite~: $\frac{\pi^2}{6}$.

\textbf{16.} Que les termes tendent vers $0$ est \emph{nécessaire}
pour que les sommes se stabilisent, mais ne décide de rien~: les
termes harmoniques $\frac1n \to 0$ et pourtant les sommes explosent~;
les termes $\frac{1}{n^2} \to 0$ et les sommes convergent. Toute la
question est de savoir à quelle \emph{vitesse} les termes meurent
--- c'est la théorie des séries, construite dans les volumes
universitaires.

\textbf{17.} $a_1 = \sqrt2 \approx 1.41421$, $b_1 = 1.5$~;
$a_2 \approx 1.45648$, $b_2 \approx 1.45711$~: deux itérations
suffisent déjà à faire coïncider trois décimales --- une vitesse
sidérante.

\textbf{18.} $b_{n+1} - a_{n+1} = \frac{a_n + b_n}{2} -
\sqrt{a_n b_n} = \frac{(\sqrt{b_n} - \sqrt{a_n})^2}{2} \geq 0$~:
les \omterm{def:g11:stat:mean}{moyennes} restent ordonnées. $(a_n)$ est \omterm{def:g12:seq:monotonic}{croissante}~:
$a_{n+1} = \sqrt{a_n b_n} \geq \sqrt{a_n \cdot a_n} = a_n$~; et
$(b_n)$ est \omterm{def:g10:functions:variations}{décroissante} par un argument symétrique.

\textbf{19.} $\frac{b_{n+1} - a_{n+1}}{b_n - a_n} =
\frac{(\sqrt{b_n} - \sqrt{a_n})^2}
{2(\sqrt{b_n} - \sqrt{a_n})(\sqrt{b_n} + \sqrt{a_n})}
= \frac{\sqrt{b_n} - \sqrt{a_n}}{2(\sqrt{b_n} + \sqrt{a_n})}
\leq \frac12$~: l'écart est au moins divisé par deux, donc
$b_n - a_n \to 0$~; avec la question 18, les \omterm{def:g12:seq:sequence}{suites} sont adjacentes
et partagent une limite $M(1, 2)$.

\textbf{20.} La troisième itération donne
$a_3 \approx b_3 \approx 1.456791$~: $M(1, 2) \approx 1.456791$ en
trois tours de manivelle (l'écart est à peu près \emph{élevé au
carré}, comme chez Héron). Classement des vitesses rencontrées dans
ce problème, de la plus lente à la plus rapide~: les sommes
harmoniques (divergence glaciale), les sommes \omterm{def:g12:seq:arith-geom}{géométriques} (erreur
divisée par deux à chaque pas), Héron et la \omterm{def:g11:stat:mean}{moyenne}
\omterm{def:g12:seq:arith-geom}{arithmético-géométrique} (erreur élevée au carré à chaque pas) ---
et c'est la vitesse surnaturelle de cette dernière qui a dit à Gauss
qu'il venait de tomber sur un nouveau filon d'analyse.
\end{solution}
